\documentclass[../readings.tex]{subfiles}

\begin{document}
\subsection{Root Finding}
\label{subsec:root-finding}

Iterative methods for solving $f(x) = 0$ for continuous $f : \fR \to \fR$.

\addDef{Stopping Tests for Root Finding}{%
Common stopping tests are
\begin{align*}
    |f(x_k)|\leq \texttt{tol}_f,
    \qquad
    |x_{k+1}-x_k|\leq \texttt{tol}_x(1+|x_k|).
\end{align*}
A small residual alone can be misleading when $f'(x^*)$ is small. A small step alone can be misleading when the iteration stagnates.
}
\label{def:root-stopping-tests}

\subsubsection{Bracketing Methods}

Guaranteed convergence to a root within a known interval.

\addThm{Bisection Method}{%
Given $f \in C[a, b]$ with $f(a)f(b) < 0$, a root $x^* \in (a, b)$ exists.
The method repeatedly halves the interval:
\begin{enumerate}
    \item $m = (a + b)/2$
    \item If $f(a)f(m) < 0$, new interval is $[a, m]$, else $[m, b]$.
\end{enumerate}
\textbf{Error Bound:}\enspace After $n$ steps, $|x_n - x^*| \leq (b - a)/2^n$.
}
\label{thm:bisection}

\addExcr{%
\begin{enumerate}
    \item Use the intermediate value theorem to prove that the initial bracket contains a root.
    \item Show that the bisection update preserves the sign-change bracket.
    \item Prove that the bracket length after $n$ steps is $(b-a)/2^n$.
    \item Derive the error bound in Thm~\ref{thm:bisection}.
\end{enumerate}
}
\label{ex:bisection}

\addRemark{%
\textbf{(Bisection trade-off)}\enspace Bisection is robust because it preserves a bracket, but its convergence is only linear. Each step gains exactly one bit of accuracy.
}

\addRemark{%
\textbf{(Bracketing default)}\enspace If a sign-changing interval is known and function evaluations are cheap, use a bracketed method. It gives a certificate that a root remains inside the current interval.
}

\subsubsection{Open Methods}

Faster local convergence, but sensitive to the initial guess $x_0$ and may diverge.

\addDef{Fixed-Point Iteration}{%
A root problem can be rewritten as $x=g(x)$ and iterated by
\begin{align*}
    x_{k+1}=g(x_k).
\end{align*}
If $|g'(x^*)|<1$, the fixed point is locally attracting.
}
\label{def:fixed-point-iteration}

\addExcr{%
\begin{enumerate}
    \item Use the mean value theorem to show $|x_{k+1}-x^*|\leq L|x_k-x^*|$ if $|g'(x)|\leq L<1$ near $x^*$.
    \item Rewrite Newton's method as a fixed-point iteration.
    \item Compute $g'(x^*)$ for Newton's fixed-point map at a simple root.
\end{enumerate}
}
\label{ex:fixed-point-iteration}

\addDef{Newton's Method}{%
Linearizes $f$ at the current iterate:
\begin{align*}
    x_{k+1} = x_k - \frac{f(x_k)}{f'(x_k)}.
\end{align*}
}
\label{def:newton-method}

\addThm{Quadratic Convergence}{%
If $f''(x)$ is continuous near a simple root $x^*$ ($f'(x^*) \neq 0$), Newton's method converges quadratically:
\begin{align*}
    |x_{k+1} - x^*| \approx \frac{|f''(x^*)|}{2|f'(x^*)|} |x_k - x^*|^2.
\end{align*}
}
\label{thm:newton-quadratic}

\addExcr{%
\begin{enumerate}
    \item Expand $f(x^*)$ about $x_k$ using a second-order Taylor formula.
    \item Divide by $f'(x_k)$ and isolate the Newton correction $f(x_k)/f'(x_k)$.
    \item Substitute $x_{k+1}=x_k-f(x_k)/f'(x_k)$.
    \item Show that $|x_{k+1}-x^*| \leq C|x_k-x^*|^2$ once $x_k$ is close enough to a simple root.
    \item Explain why the assumption $f'(x^*)\neq 0$ is essential.
\end{enumerate}
}
\label{ex:newton-quadratic}

\addRemark{%
\textbf{(Newton failure modes)}\enspace
1. $x_0$ too far from $x^*$.
2. $f'(x_k) \approx 0$ (zero derivative).
3. Multiple roots ($f'(x^*) = 0$): convergence degrades to linear with rate $1/2$.
}

\addThm{Newton on a Multiple Root}{%
If $f(x)=(x-x^*)^m q(x)$ with $q(x^*)\neq 0$ and $m>1$, then Newton's method converges locally linearly with asymptotic factor
\begin{align*}
    1-\frac{1}{m}.
\end{align*}
}
\label{thm:newton-multiple-root}

\addExcr{%
\begin{enumerate}
    \item Compute $f'(x)$ for $f(x)=(x-x^*)^m q(x)$.
    \item Substitute $f$ and $f'$ into the Newton update.
    \item Keep the leading term as $x_k\to x^*$.
    \item Derive the factor in Thm~\ref{thm:newton-multiple-root}.
\end{enumerate}
}
\label{ex:newton-multiple-root}

\addDef{Secant Method}{%
Approximates $f'$ via finite differences of the last two iterates:
\begin{align*}
    x_{k+1} = x_k - f(x_k)\frac{x_k - x_{k-1}}{f(x_k) - f(x_{k-1})}.
\end{align*}
Requires two initial points ($x_0, x_1$).
}
\label{def:secant-method}

\addThm{Superlinear Convergence}{%
The secant method converges with order $\phi = (1 + \sqrt{5})/2 \approx 1.618$ (the golden ratio).
\textbf{Efficiency:}\enspace Faster than bisection, slower than Newton, but requires only 1 function evaluation per step (Newton requires 2: $f$ and $f'$).
}
\label{thm:secant-superlinear}

\addExcr{%
\begin{enumerate}
    \item Derive the secant update by replacing $f'(x_k)$ in Newton's method with a finite difference through $(x_{k-1},f(x_{k-1}))$ and $(x_k,f(x_k))$.
    \item Compare the information used by Newton and secant.
    \item Numerically estimate the convergence order for $f(x)=x^2-2$.
\end{enumerate}
}
\label{ex:secant-method}

\subsubsection{Brent's Method}

A standard robust algorithm for one-dimensional root-finding.

\addThm{The Hybrid Approach}{%
Brent's method combines:
\begin{enumerate}
    \item \textbf{Bisection:}\enspace For guaranteed global convergence.
    \item \textbf{Secant:}\enspace For fast local convergence.
    \item \textbf{IQI:}\enspace Inverse Quadratic Interpolation using 3 points.
\end{enumerate}
It uses fast open steps by default, falling back to bisection only if the step is outside the bracket or not improving quickly enough.
}
\label{thm:brent-hybrid}

\addRemark{%
\textbf{Implementation:}\enspace Used in \texttt{scipy.optimize.brentq}.
}

\addDef{Method Selection}{%
\begin{center}
\begin{tabular}{ll}
\textbf{Available information} & \textbf{Typical method} \\ \hline
Sign-changing bracket & Bisection or Brent \\
Good initial guess and derivative & Newton \\
Good initial guesses, no derivative & Secant \\
Need robust scalar default & Brent \\
\end{tabular}
\end{center}
}
\label{def:root-method-selection}

\subsubsection{Exercises}

\addExcr{%
\begin{enumerate}
    \item Bisection on $f(x) = x^3 - 2$ on $[1, 2]$ for 4 steps. Find the guaranteed error using Thm~\ref{thm:bisection}.
    \item Newton on $f(x) = x^2 - 2$ from $x_0 = 1$. Verify the doubling of digits using Thm~\ref{thm:newton-quadratic}.
    \item Prove that for $f(x) = x^2$, Newton converges to $0$ with linear rate $1/2$ using Thm~\ref{thm:newton-multiple-root}.
    \item Find the root of $e^x = 3x$ near $x=1$ via Secant method. Compare evaluations vs.\ Newton.
    \item Construct an example where $|f(x_k)|$ is small but $|x_k-x^*|$ is not small. Relate it to Def~\ref{def:root-stopping-tests}.
\end{enumerate}
}

\end{document}
